\chapter{HASIL DAN PEMBAHASAN}

\section{Data dan Pelatihan Model Final}
Model final (\kode{hmf-aug-v2}) dilatih dengan 35.766 lagu (Tabel~\ref{tab:data-final}). Pelatihan berjalan 30 epoch di GPU Google Colab dan dilanjutkan dari epoch 8 setelah sesi terputus. Epoch 27 dipilih karena loss validasinya terendah, yaitu 0,184 (Gambar~\ref{fig:loss}).

\begin{table}[htbp]
\centering
\singlespacing\small
\caption{Pembagian data model final.}
\label{tab:data-final}
\begin{tabular}{lrrrr}
\toprule
Bagian & Lagu & Manusia & AI & Segmen 5 detik\\
\midrule
Train & 28.587 (79,9\%) & 8.198 & 20.389 & 508.188\textsuperscript{a}\\
Validation & 3.268 (9,1\%) & 2.158 & 1.110 & 19.237\\
Test & 3.911 (10,9\%) & 1.330 & 2.581 & 23.083\\
\midrule
Total & 35.766 & 11.686 & 24.080 & \\
\bottomrule
\multicolumn{5}{l}{\footnotesize \textsuperscript{a} Termasuk 3 versi augmentasi codec per lagu.}
\end{tabular}
\end{table}

\begin{figure}[htbp]
\centering
\begin{tikzpicture}
\begin{axis}[
  width=12.5cm, height=6.5cm,
  xmin=1, xmax=30, ymin=0, ymax=0.45,
  xtick={1,5,10,15,20,25,30},
  ytick={0,0.1,0.2,0.3,0.4},
  yticklabel style={/pgf/number format/.cd, use comma, fixed, precision=1},
  xlabel={Epoch}, ylabel={Loss},
  axis lines*=left, grid=major, grid style={black!8},
  legend style={draw=none, font=\small, at={(0.98,0.95)}, anchor=north east},
  tick label style={font=\small}, label style={font=\small}]
\addplot[thick, color=black!55] coordinates {(1,0.4317) (2,0.2678) (3,0.2104) (4,0.1775) (5,0.1576) (6,0.1430) (7,0.1326) (8,0.1226) (9,0.1182) (10,0.1111) (11,0.1058) (12,0.1032) (13,0.0980) (14,0.0951) (15,0.0923) (16,0.0898) (17,0.0865) (18,0.0852) (19,0.0822) (20,0.0813) (21,0.0791) (22,0.0775) (23,0.0757) (24,0.0744) (25,0.0739) (26,0.0718) (27,0.0699) (28,0.0687) (29,0.0688) (30,0.0666)};
\addlegendentry{Loss train}
\addplot[thick, color=aksen] coordinates {(1,0.3987) (2,0.3453) (3,0.3044) (4,0.3035) (5,0.2668) (6,0.2742) (7,0.2534) (8,0.2428) (9,0.2457) (10,0.2342) (11,0.2394) (12,0.2585) (13,0.2457) (14,0.2504) (15,0.2188) (16,0.2251) (17,0.2197) (18,0.2459) (19,0.2117) (20,0.2162) (21,0.2266) (22,0.2061) (23,0.2184) (24,0.1986) (25,0.1866) (26,0.2011) (27,0.1839) (28,0.2055) (29,0.1896) (30,0.1973)};
\addlegendentry{Loss validasi}
\addplot[only marks, mark=*, mark size=2.6pt, color=aksen, forget plot] coordinates {(27,0.1839)};
\node[font=\footnotesize, anchor=south] at (axis cs:27,0.192) {epoch terpilih};
\end{axis}
\end{tikzpicture}
\caption{Loss train dan validasi per epoch model final; titik menandai epoch 27 yang dipilih.}
\label{fig:loss}
\end{figure}

\section{Hasil Hold-Out Test}
Pada 3.911 lagu test, model final mencapai akurasi 86,5\% dan ROC AUC 0,966 (Tabel~\ref{tab:test}). Salah tuduh jarang terjadi: hanya 47 dari 1.330 lagu manusia (3,5\%) divonis AI (Tabel~\ref{tab:confusion}). Hampir semua kesalahan adalah lagu AI yang lolos sebagai manusia.

\begin{table}[htbp]
\centering
\singlespacing\small
\caption{Hasil hold-out test model final.}
\label{tab:test}
\begin{tabular}{lr}
\toprule
Metrik & Nilai\\
\midrule
Accuracy & 86,5\%\\
Balanced accuracy & 88,9\%\\
Macro F1 & 0,859\\
ROC AUC & 0,966\\
Recall manusia (lagu manusia benar dikenali) & 96,5\% (1.283 dari 1.330)\\
Recall AI (lagu AI terdeteksi) & 81,4\% (2.101 dari 2.581)\\
Presisi AI (vonis AI yang memang AI) & 97,8\%\\
\bottomrule
\end{tabular}
\end{table}

\begin{table}[htbp]
\centering
\singlespacing\small
\caption{Confusion matrix model final pada data test.}
\label{tab:confusion}
\begin{tabular}{lrr}
\toprule
Asal lagu sebenarnya & Divonis human & Divonis AI\\
\midrule
Manusia & 1.283 & 47\\
AI & 480 & 2.101\\
\bottomrule
\end{tabular}
\end{table}

\section{Recall per Generator}
Semua generator terdeteksi minimal 91\%, kecuali udio-30s yang hanya 35,1\% (Gambar~\ref{fig:per-generator}). Tanpa udio-30s, recall AI naik menjadi 98,5\% (1.857 dari 1.886 lagu). Generator udio-30s, chirp-v3, dan chirp-v2-xxl-alpha sama sekali tidak ada di data train dan hanya muncul di validation dan test. Dua versi Suno itu tetap terdeteksi sekitar 99\%, tetapi udio-30s tidak. Artinya, model dapat mengenali versi baru dari generator yang sudah dikenalnya, tetapi tidak selalu. Inilah batasan generalisasi yang paling penting.

\begin{figure}[htbp]
\centering
\begin{tikzpicture}
\begin{axis}[
  xbar, bar width=9pt, bar shift=0pt, y=0.72cm,
  xmin=0, xmax=118, xtick={0,25,50,75,100},
  ytick={1,...,11},
  yticklabels={udio-30s\,*, stable\_audio, jasco, udio-120s, magenta, audioldm, chirp-v2-xxl-alpha\,*, riffusion, chirp-v3\,*, chirp-v3.5, musicgen},
  xlabel={Recall pada data test (\%)},
  axis x line*=bottom, axis y line*=left,
  enlarge y limits=0.05, width=10.4cm,
  tick label style={font=\small}, label style={font=\small},
  clip=false]
\addplot[fill=black!35, draw=none] coordinates {
  (91.3,2) (93.4,3) (97.9,4) (98.0,5) (98.0,6) (98.9,7) (98.9,8) (99.8,9) (100,10) (100,11)};
\addplot[fill=aksen, draw=none] coordinates {(35.1,1)};
\foreach \x/\y/\t in {35.1/1/{35,1\% (244/695)}, 91.3/2/{91,3\% (95/104)}, 93.4/3/{93,4\% (85/91)},
  97.9/4/{97,9\% (238/243)}, 98.0/5/{98,0\% (96/98)}, 98.0/6/{98,0\% (97/99)},
  98.9/7/{98,9\% (266/269)}, 98.9/8/{98,9\% (89/90)}, 99.8/9/{99,8\% (558/559)},
  100/10/{100\% (234/234)}, 100/11/{100\% (99/99)}}
  {\edef\temp{\noexpand\node[anchor=west, font=\noexpand\footnotesize, inner sep=2pt] at (axis cs:\x,\y) {\unexpanded\expandafter{\t}};}\temp}
\end{axis}
\end{tikzpicture}
\caption{Recall per generator pada data test model final; tanda * berarti generator tidak ada di data train.}
\label{fig:per-generator}
\end{figure}

\section{Uji Robustness terhadap Kompresi Ulang}
Uji \textit{robustness} memakai 200 lagu test (100 per kelas) yang dikompres ulang dengan empat format, lalu membandingkan vonis model final sebelum dan sesudah kompresi (Tabel~\ref{tab:robust}).

\begin{table}[htbp]
\centering
\singlespacing\small
\caption{Hasil uji robustness model final.}
\label{tab:robust}
\begin{tabular}{lccc}
\toprule
Kondisi & Akurasi & Flip rate & Mean $|\Delta|$ skor\\
\midrule
Asli (clean) & 88,5\% & -- & --\\
MP3 32 kbps mono 16 kHz & 86,5\% & 5,0\% & 0,055\\
MP3 128 kbps & 88,0\% & 0,5\% & 0,012\\
AAC 128 kbps & 88,0\% & 0,5\% & 0,016\\
Opus 64 kbps & 86,0\% & 3,5\% & 0,035\\
\bottomrule
\end{tabular}
\end{table}

Akurasi model final bertahan di 86--88\% pada semua kondisi, dengan ROC AUC 0,959 pada file asli. Kondisi terberat, MP3 32 kbps mono 16 kHz, hanya menurunkan akurasi 2 poin dan membalik 5\% vonis; rata-rata perubahan skor hanya 0,055. Pada MP3 dan AAC 128 kbps, skor hampir tidak bergeser. Artinya, vonis model ditentukan oleh isi musik, bukan oleh format atau kualitas file.

\section{Studi Ablasi}
Studi ablasi menambahkan perbaikan satu per satu untuk melihat peran masing-masing. Percobaan dilakukan dalam skala kecil, yaitu 1.000 lagu dengan 2 epoch di CPU dan 20 lagu test per kelas (Tabel~\ref{tab:ablasi}).

\begin{table}[htbp]
\centering
\singlespacing\small
\caption{Studi ablasi skala kecil.}
\label{tab:ablasi}
\begin{tabularx}{\textwidth}{Y L{2.5cm} L{1.4cm} L{2.0cm} L{1.6cm}}
\toprule
Model & Skor manusia: asli $\rightarrow$ MP3 32k & Flip MP3 32k & $|\Delta|$ kondisi terburuk & Akurasi asli\\
\midrule
Tanpa augmentasi & 0,00 $\rightarrow$ 1,00 & 50\% & 0,501 & 100\%\\
Augmentasi tanpa \kode{--fmax} & 0,04 $\rightarrow$ 0,91 & 50\% & 0,442 & 100\%\\
Augmentasi dan \kode{--fmax 7000} & 0,33 $\rightarrow$ 0,33 & 0\% & 0,008 & 77,5\%\\
\bottomrule
\end{tabularx}
\end{table}

Augmentasi saja belum cukup: vonis tetap berbalik pada 50\% lagu. Baru setelah spektrogram dibatasi sampai 7 kHz, skor tidak lagi berubah oleh codec. Akurasi 100\% pada dua baris pertama berasal dari petunjuk kualitas file, sehingga langsung runtuh begitu lagu manusia dikompres. Angka akurasi pada tabel ini berasal dari model kecil 2 epoch; yang dibaca adalah polanya.

\section{Hasil Pengujian Perangkat Lunak}\label{sec:hasil-blackbox}
Seluruh 42 tes otomatis pytest lulus. Hasil \textit{black-box testing} studio web disajikan pada Tabel~\ref{tab:blackbox}. Skenario ini menguji perilaku aplikasi, bukan akurasi model.

{\singlespacing\small
\begin{longtable}{L{0.5cm} L{2.3cm} L{3.2cm} L{4.4cm} L{1.2cm}}
\caption{Hasil black-box testing studio web.}\label{tab:blackbox}\\
\toprule
No & Skenario & Masukan & Hasil yang diharapkan & Hasil\\
\midrule
\endfirsthead
\toprule
No & Skenario & Masukan & Hasil yang diharapkan & Hasil\\
\midrule
\endhead
\bottomrule
\endfoot
1 & Membuka studio & Buka \url{http://127.0.0.1:8765} & Halaman tampil dengan nama model aktif & Sesuai\\
2 & Analisis lagu & MP3 2 menit & Vonis, skor AI, spektrogram, dan 24 batang skor segmen tampil & Sesuai\\
3 & Kecocokan skor & Lagu yang sama & Persentase di layar sama dengan \kode{ai\_score} di JSON & Sesuai\\
4 & File bukan audio & File teks & Pesan ``Audio gagal dibaca'' & Sesuai\\
5 & Port sudah dipakai & Studio kedua di port yang sama & Server menolak berjalan dan menyarankan port lain & Sesuai\\
6 & Simpan umpan balik & Tandai vonis salah, pilih sumber label, simpan audio & Umpan balik tersimpan; panel uji lapangan bertambah & Sesuai\\
7 & Audio berbeda & File lain dikirim sebagai audio umpan balik & Ditolak dengan pesan ``File berbeda dari lagu yang dianalisis'' & Sesuai\\
8 & Tampilan ponsel & Lebar layar 375 piksel & Semua bagian tersusun satu kolom tanpa geser ke samping & Sesuai\\
\end{longtable}}

\section{Pembahasan dan Keterbatasan}
Hasil di atas menjawab rumusan masalah. Detektor tidak bergantung pada format file: vonisnya stabil pada empat kondisi kompresi ulang. Akurasi test 86,5\% dengan ROC AUC 0,966 mencerminkan kemampuan yang sebenarnya karena data latih sudah diaugmentasi dan dibatasi sampai 7 kHz, sehingga petunjuk codec tidak bisa dipakai. Model juga berhati-hati: hanya 3,5\% lagu manusia yang salah dituduh AI. Hasil deteksi disajikan secara transparan melalui skor per segmen, spektrogram, dan panel model.

Keterbatasan yang perlu diperhatikan:
\begin{itemize}
\item Model hanya mengenali pola yang ada di data latih. Generator yang belum pernah dilihat, seperti udio-30s, dapat lolos sebagai manusia.
\item Ambang 0,5 belum dikalibrasi; skor 51\% tidak berarti peluang 51\% bahwa lagu itu buatan AI.
\item Spektrogram hanya sampai 7 kHz, sehingga artefak AI di frekuensi tinggi tidak terlihat oleh model.
\item Lagu campuran manusia dan AI belum menjadi kelas tersendiri.
\item Sebagian besar kesalahan adalah lagu AI yang lolos sebagai manusia (480 dari 2.581 lagu AI), terutama dari udio-30s.
\end{itemize}

Hasil studio tidak boleh dipakai sebagai satu-satunya dasar untuk menuduh seseorang memakai AI. Skor sebaiknya dilaporkan sebagai perkiraan bersama nama model dan grafik per segmen, serta digabungkan dengan bukti lain seperti file proyek atau rekaman sesi.
